\documentclass{article}
\usepackage{amsmath, amssymb, geometry, xcolor}
\geometry{letterpaper, margin=1in}

\title{Homework 2: Logic (Tautologies, Contradictions, and Logical Equivalence)}
\author{Tashfeen Omran}
\date{September 2026}

\begin{document}
\maketitle

\section*{Part 1: Comprehensive Introduction, Context, and Prerequisites}

\subsection*{Core Concepts}
In the study of formal logic, we analyze the structure of statements and how they combine to form complex arguments. The foundation of this system rests on \textbf{propositions} (or statements)---declarative sentences that are unambiguously either true (T) or false (F). When we connect these statements using logical connectives such as AND ($\land$), OR ($\lor$), NOT ($\sim$), IMPLIES ($\Rightarrow$), and IF AND ONLY IF ($\Leftrightarrow$), we create \textbf{compound statements}.

Three core concepts emerge when we evaluate the truth values of these compound statements across all possible scenarios (using truth tables):
\begin{itemize}
    \item \textbf{Tautology:} A compound statement that is \textit{always true}, regardless of the truth values of its component statements. It represents an undeniable logical certainty.
    \item \textbf{Contradiction:} A compound statement that is \textit{always false}, representing a logical impossibility.
    \item \textbf{Logical Equivalence ($\equiv$):} Two statements are logically equivalent if they share the exact same truth values in every possible scenario. If $P \equiv Q$, you can replace $P$ with $Q$ in any logical argument without changing the outcome.
\end{itemize}

\subsection*{Intuition and Derivation}
Why do we rely on truth tables? A truth table is an exhaustive algorithmic check. For a compound statement with $n$ distinct variables, there are exactly $2^n$ possible combinations of true/false states. By iterating through every single permutation, a truth table acts as a mathematical proof by exhaustion. 

The equivalence $P \Rightarrow Q \equiv \sim P \lor Q$ is notoriously counter-intuitive for beginners, but derives perfectly from how we define a "broken promise." If I say, "If you get an A, I will buy you a car" ($P \Rightarrow Q$), I have only lied to you if you get an A ($P$ is true) AND I do not buy the car ($Q$ is false). In every other scenario (you didn't get an A, or I bought the car anyway), my initial statement was not violated. Thus, the implication is true as long as $P$ is false ($\sim P$) OR $Q$ is true ($Q$). 

\subsection*{Historical Context and Motivation}
The formalization of logical structures traces back to ancient Greek philosophy with Aristotle's syllogisms. However, the modern algebraic approach to logic was revolutionized by George Boole in the 1840s. Boole sought to translate human reasoning and philosophical arguments into a rigorous, mathematical calculus where logical deductions could be computed like algebraic equations. 

This motivation was profound: mathematicians in the 19th century were encountering paradoxes in set theory and analysis. They needed an absolute, infallible foundation for mathematical proofs. By establishing rules for tautologies and equivalences, Boole (and later logicians like Frege and Russell) provided the rigorous framework needed to ensure that mathematical theorems were proven without logical fallacies.

\subsection*{Key Formulas and Equivalences}
\begin{itemize}
    \item \textbf{Implication as Disjunction:} $P \Rightarrow Q \equiv \sim P \lor Q$
    \item \textbf{Negation of an Implication:} $\sim(P \Rightarrow Q) \equiv P \land \sim Q$
    \item \textbf{Contrapositive:} $P \Rightarrow Q \equiv \sim Q \Rightarrow \sim P$
    \item \textbf{De Morgan's Laws:} 
    \begin{align*}
        \sim(P \land Q) &\equiv \sim P \lor \sim Q \\
        \sim(P \lor Q) &\equiv \sim P \land \sim Q
    \end{align*}
    \item \textbf{Distributive Laws:}
    \begin{align*}
        P \lor (Q \land R) &\equiv (P \lor Q) \land (P \lor R) \\
        P \land (Q \lor R) &\equiv (P \land Q) \lor (P \land R)
    \end{align*}
\end{itemize}

\subsection*{Prerequisites}
To master this topic, you must be comfortable with basic propositional variables, the five standard logical connectives ($\sim, \lor, \land, \Rightarrow, \Leftrightarrow$), and the combinatorial setup of truth tables (knowing how to alternate T and F systematically to cover all $2^n$ rows).

\newpage
\section*{Part 2: Detailed Homework Solutions}

\subsection*{Section 2.7 Exercises}

\textbf{Problem 2.46} For statements $P$ and $Q$, show that $P \Rightarrow (P \lor Q)$ is a tautology.
\\[0.5em]
\textbf{Solution:} We construct a truth table for $P \Rightarrow (P \lor Q)$.
\[
\begin{array}{|c|c|c|c|}
\hline
P & Q & P \lor Q & P \Rightarrow (P \lor Q) \\
\hline
\text{T} & \text{T} & \text{T} & \text{T} \\
\text{T} & \text{F} & \text{T} & \text{T} \\
\text{F} & \text{T} & \text{T} & \text{T} \\
\text{F} & \text{F} & \text{F} & \text{T} \\
\hline
\end{array}
\]
Since all entries in the final column are True (T), the statement $P \Rightarrow (P \lor Q)$ is a tautology.

\vspace{1em}
\textbf{Problem 2.47} For statements $P$ and $Q$, show that $(P \land (\sim Q)) \land (P \land Q)$ and $(P \Rightarrow \sim Q) \land (P \land Q)$ are contradictions.
\\[0.5em]
\textbf{Solution:}
First, we analyze $(P \land \sim Q) \land (P \land Q)$:
\[
\begin{array}{|c|c|c|c|c|c|}
\hline
P & Q & \sim Q & P \land \sim Q & P \land Q & (P \land \sim Q) \land (P \land Q) \\
\hline
\text{T} & \text{T} & \text{F} & \text{F} & \text{T} & \text{F} \\
\text{T} & \text{F} & \text{T} & \text{T} & \text{F} & \text{F} \\
\text{F} & \text{T} & \text{F} & \text{F} & \text{F} & \text{F} \\
\text{F} & \text{F} & \text{T} & \text{F} & \text{F} & \text{F} \\
\hline
\end{array}
\]
Since the final column is all False, it is a contradiction.

Next, we analyze $(P \Rightarrow \sim Q) \land (P \land Q)$:
\[
\begin{array}{|c|c|c|c|c|c|}
\hline
P & Q & \sim Q & P \Rightarrow \sim Q & P \land Q & (P \Rightarrow \sim Q) \land (P \land Q) \\
\hline
\text{T} & \text{T} & \text{F} & \text{F} & \text{T} & \text{F} \\
\text{T} & \text{F} & \text{T} & \text{T} & \text{F} & \text{F} \\
\text{F} & \text{T} & \text{F} & \text{T} & \text{F} & \text{F} \\
\text{F} & \text{F} & \text{T} & \text{T} & \text{F} & \text{F} \\
\hline
\end{array}
\]
Again, the final column is exclusively False, confirming it is a contradiction.

\vspace{1em}
\textbf{Problem 2.48} For statements $P$ and $Q$, show that $((P \Rightarrow Q) \land P) \Rightarrow Q$ is a tautology. Then state this logical argument form in words.
\\[0.5em]
\textbf{Solution:}
\[
\begin{array}{|c|c|c|c|c|}
\hline
P & Q & P \Rightarrow Q & (P \Rightarrow Q) \land P & ((P \Rightarrow Q) \land P) \Rightarrow Q \\
\hline
\text{T} & \text{T} & \text{T} & \text{T} & \text{T} \\
\text{T} & \text{F} & \text{F} & \text{F} & \text{T} \\
\text{F} & \text{T} & \text{T} & \text{F} & \text{T} \\
\text{F} & \text{F} & \text{T} & \text{F} & \text{T} \\
\hline
\end{array}
\]
The final column is all True, verifying it is a tautology.
\textbf{In words:} ``If $P$ implies $Q$, and $P$ is true, then $Q$ is true.'' (This is known as \textit{modus ponens}).

\vspace{1em}
\textbf{Problem 2.49} For statements $P, Q$ and $R$, show that $((P \Rightarrow Q) \land (Q \Rightarrow R)) \Rightarrow (P \Rightarrow R)$ is a tautology.
\\[0.5em]
\textbf{Solution:} Let $\Omega$ be the statement $((P \Rightarrow Q) \land (Q \Rightarrow R)) \Rightarrow (P \Rightarrow R)$.
\[
\begin{array}{|c|c|c|c|c|c|c|c|}
\hline
P & Q & R & P \Rightarrow Q & Q \Rightarrow R & (P \Rightarrow Q) \land (Q \Rightarrow R) & P \Rightarrow R & \Omega \\
\hline
\text{T} & \text{T} & \text{T} & \text{T} & \text{T} & \text{T} & \text{T} & \text{T} \\
\text{T} & \text{T} & \text{F} & \text{T} & \text{F} & \text{F} & \text{F} & \text{T} \\
\text{T} & \text{F} & \text{T} & \text{F} & \text{T} & \text{F} & \text{T} & \text{T} \\
\text{T} & \text{F} & \text{F} & \text{F} & \text{T} & \text{F} & \text{F} & \text{T} \\
\text{F} & \text{T} & \text{T} & \text{T} & \text{T} & \text{T} & \text{T} & \text{T} \\
\text{F} & \text{T} & \text{F} & \text{T} & \text{F} & \text{F} & \text{T} & \text{T} \\
\text{F} & \text{F} & \text{T} & \text{T} & \text{T} & \text{T} & \text{T} & \text{T} \\
\text{F} & \text{F} & \text{F} & \text{T} & \text{T} & \text{T} & \text{T} & \text{T} \\
\hline
\end{array}
\]
The statement evaluates to True in all 8 possible cases; it is a tautology (the \textit{Law of Syllogism}).

\vspace{1em}
\textbf{Problem 2.50} Let $R$ and $S$ be compound statements... If $R$ is a tautology and $S$ is a contradiction, then what can be said of the following?
\\[0.5em]
\textbf{Solution:} We treat $R$ as always True (T) and $S$ as always False (F).
\begin{itemize}
    \item \textbf{(a) $R \lor S$:} Evaluates to $\text{T} \lor \text{F} \equiv \text{T}$. It is a \textbf{tautology}.
    \item \textbf{(b) $R \land S$:} Evaluates to $\text{T} \land \text{F} \equiv \text{F}$. It is a \textbf{contradiction}.
    \item \textbf{(c) $R \Rightarrow S$:} Evaluates to $\text{T} \Rightarrow \text{F} \equiv \text{F}$. It is a \textbf{contradiction}.
    \item \textbf{(d) $S \Rightarrow R$:} Evaluates to $\text{F} \Rightarrow \text{T} \equiv \text{T}$. It is a \textbf{tautology}.
\end{itemize}

\vspace{1em}
\textbf{Problem 2.51} For statements $P$ and $Q$, determine whether $(P \lor Q) \lor (Q \Rightarrow P)$ is a tautology, contradiction, or neither.
\\[0.5em]
\textbf{Solution:} Using algebraic laws of logic:
\[ (P \lor Q) \lor (Q \Rightarrow P) \equiv P \lor Q \lor (\sim Q \lor P) \]
By associativity and commutativity:
\[ \equiv (P \lor P) \lor (Q \lor \sim Q) \]
Since $P \lor P \equiv P$, and $Q \lor \sim Q \equiv \text{T}$:
\[ \equiv P \lor \text{T} \equiv \text{T} \]
The statement is a \textbf{tautology}.

\vspace{1em}
\textbf{Problem 2.52} Determine whether $((P \Rightarrow Q) \Rightarrow P) \Rightarrow (P \Rightarrow (Q \Rightarrow P))$ is a tautology, contradiction, or neither.
\\[0.5em]
\textbf{Solution:} Let's evaluate the conclusion of the main implication, $P \Rightarrow (Q \Rightarrow P)$:
\[ P \Rightarrow (Q \Rightarrow P) \equiv \sim P \lor (\sim Q \lor P) \equiv (\sim P \lor P) \lor \sim Q \equiv \text{T} \lor \sim Q \equiv \text{T} \]
Because the conclusion is a tautology (always True), the overall structure is $(\text{Something}) \Rightarrow \text{T}$, which is always True. Therefore, it is a \textbf{tautology}.

\subsection*{Section 2.8 Exercises}

\textbf{Problem 2.53} The implication $(\sim P) \Rightarrow (\sim Q)$ is the inverse of $P \Rightarrow Q$. (a) Show these are not logically equivalent. (b) Find another implication equivalent to the inverse.
\\[0.5em]
\textbf{Solution (a):}
\[
\begin{array}{|c|c|c|c|c|c|}
\hline
P & Q & \sim P & \sim Q & P \Rightarrow Q & \sim P \Rightarrow \sim Q \\
\hline
\text{T} & \text{T} & \text{F} & \text{F} & \text{T} & \text{T} \\
\text{T} & \text{F} & \text{F} & \text{T} & \text{F} & \text{T} \\
\text{F} & \text{T} & \text{T} & \text{F} & \text{T} & \text{F} \\
\text{F} & \text{F} & \text{T} & \text{T} & \text{T} & \text{T} \\
\hline
\end{array}
\]
Rows 2 and 3 differ. Therefore, they are \textbf{not logically equivalent}.
\\
\textbf{Solution (b):} The converse $Q \Rightarrow P$ is logically equivalent to the inverse $\sim P \Rightarrow \sim Q$. 
Proof by contrapositive: the contrapositive of $Q \Rightarrow P$ is $\sim P \Rightarrow \sim Q$.

\vspace{1em}
\textbf{Problem 2.54} (a) Is $\sim(P \lor Q)$ logically equivalent to $(\sim P) \lor (\sim Q)$? (b) What can you say about $\sim(P \lor Q) \Leftrightarrow ((\sim P) \lor (\sim Q))$?
\\[0.5em]
\textbf{Solution (a):} \textbf{No.} By De Morgan's Law, $\sim(P \lor Q) \equiv \sim P \land \sim Q$, which differs from $\sim P \lor \sim Q$. For example, if $P$ is F and $Q$ is T, $\sim(P \lor Q)$ is F, but $\sim P \lor \sim Q$ is T.
\\
\textbf{Solution (b):} Because the two sides of the biconditional are not logically equivalent, the biconditional is \textbf{not a tautology}. It will evaluate to False whenever $P$ and $Q$ possess differing truth values, and will only evaluate to True if $P$ and $Q$ are both True or both False.

\vspace{1em}
\textbf{Problem 2.55} Use a truth table to show equivalence:
(a) $(P \land Q) \Leftrightarrow P$ and $P \Rightarrow Q$.
(b) $P \Rightarrow (Q \lor R)$ and $(P \land \sim Q) \Rightarrow R$.
\\[0.5em]
\textbf{Solution (a):}
\[
\begin{array}{|c|c|c|c|c|}
\hline
P & Q & P \land Q & (P \land Q) \Leftrightarrow P & P \Rightarrow Q \\
\hline
\text{T} & \text{T} & \text{T} & \text{T} & \text{T} \\
\text{T} & \text{F} & \text{F} & \text{F} & \text{F} \\
\text{F} & \text{T} & \text{F} & \text{T} & \text{T} \\
\text{F} & \text{F} & \text{F} & \text{T} & \text{T} \\
\hline
\end{array}
\]
The last two columns match exactly, proving logical equivalence.

\textbf{Solution (b):}
\[
\begin{array}{|c|c|c|c|c|c|c|c|}
\hline
P & Q & R & Q \lor R & P \Rightarrow (Q \lor R) & P \land \sim Q & (P \land \sim Q) \Rightarrow R \\
\hline
\text{T} & \text{T} & \text{T} & \text{T} & \text{T} & \text{F} & \text{T} \\
\text{T} & \text{T} & \text{F} & \text{T} & \text{T} & \text{F} & \text{T} \\
\text{T} & \text{F} & \text{T} & \text{T} & \text{T} & \text{T} & \text{T} \\
\text{T} & \text{F} & \text{F} & \text{F} & \text{F} & \text{T} & \text{F} \\
\text{F} & \text{T} & \text{T} & \text{T} & \text{T} & \text{F} & \text{T} \\
\text{F} & \text{T} & \text{F} & \text{T} & \text{T} & \text{F} & \text{T} \\
\text{F} & \text{F} & \text{T} & \text{T} & \text{T} & \text{F} & \text{T} \\
\text{F} & \text{F} & \text{F} & \text{F} & \text{T} & \text{F} & \text{T} \\
\hline
\end{array}
\]
Columns 5 and 7 match exactly. They are logically equivalent.

\vspace{1em}
\textbf{Problem 2.56} Show that $(\sim Q) \Rightarrow (P \land \sim P)$ and $Q$ are logically equivalent.
\\[0.5em]
\textbf{Solution:} Notice that $P \land \sim P$ is a contradiction (always False).
So, the implication simplifies to $\sim Q \Rightarrow \text{F}$. 
If $Q$ is True, $\sim Q$ is False, making $\text{F} \Rightarrow \text{F}$, which is True.
If $Q$ is False, $\sim Q$ is True, making $\text{T} \Rightarrow \text{F}$, which is False.
In both scenarios, the truth value matches $Q$. Thus, $(\sim Q) \Rightarrow (P \land \sim P) \equiv Q$.

\vspace{1em}
\textbf{Problem 2.57} Show that $(P \lor Q) \Rightarrow R$ and $(P \Rightarrow R) \land (Q \Rightarrow R)$ are equivalent.
\\[0.5em]
\textbf{Solution:} By replacing implications with disjunctions:
\[ (P \lor Q) \Rightarrow R \equiv \sim(P \lor Q) \lor R \]
Apply De Morgan's Law:
\[ \equiv (\sim P \land \sim Q) \lor R \]
Apply the Distributive Law:
\[ \equiv (\sim P \lor R) \land (\sim Q \lor R) \]
Convert back to implications:
\[ \equiv (P \Rightarrow R) \land (Q \Rightarrow R) \]
This algebraically verifies the equivalence.

\vspace{1em}
\textbf{Problem 2.58} If $S$ and $T$ are not logically equivalent, what can we conclude?
\\[0.5em]
\textbf{Solution:} We can definitively conclude that there exists at least one specific assignment of truth values to the component statements $P, Q,$ and $R$ that causes $S$ and $T$ to yield different truth values (i.e., their truth tables differ in at least one row).

\vspace{1em}
\textbf{Problem 2.59} Five statements $S_1, \dots, S_5$ made from $P$ and $Q$ have identical 1st and 4th rows. Show at least two are equivalent.
\\[0.5em]
\textbf{Solution:} A truth table with two variables ($P, Q$) has exactly $2^2 = 4$ rows. We are given that rows 1 and 4 are fixed/identical across all five statements. The only rows that can vary are rows 2 and 3. Each row can be T or F, meaning there are $2 \times 2 = 4$ possible distinct T/F combinations for rows 2 and 3. Because we have $5$ statements but only $4$ possible truth assignments for the remaining rows, the \textbf{Pigeonhole Principle} dictates that at least two statements must possess the exact same combination for rows 2 and 3. Since rows 1 and 4 were already identical, these two statements share completely identical truth tables, rendering them logically equivalent.

\subsection*{Section 2.9 Exercises}

\textbf{Problem 2.60} Verify (a) $P \lor (Q \land R) \equiv (P \lor Q) \land (P \lor R)$ and (b) $\sim(P \lor Q) \equiv \sim P \land \sim Q$.
\\[0.5em]
\textbf{Solution (a):} We create an 8-row truth table. Let $S = (P \lor Q) \land (P \lor R)$.
\[
\begin{array}{|c|c|c|c|c|c|c|c|}
\hline
P & Q & R & Q \land R & P \lor (Q \land R) & P \lor Q & P \lor R & S \\
\hline
\text{T} & \text{T} & \text{T} & \text{T} & \text{T} & \text{T} & \text{T} & \text{T} \\
\text{T} & \text{T} & \text{F} & \text{F} & \text{T} & \text{T} & \text{T} & \text{T} \\
\text{T} & \text{F} & \text{T} & \text{F} & \text{T} & \text{T} & \text{T} & \text{T} \\
\text{T} & \text{F} & \text{F} & \text{F} & \text{T} & \text{T} & \text{T} & \text{T} \\
\text{F} & \text{T} & \text{T} & \text{T} & \text{T} & \text{T} & \text{T} & \text{T} \\
\text{F} & \text{T} & \text{F} & \text{F} & \text{F} & \text{T} & \text{F} & \text{F} \\
\text{F} & \text{F} & \text{T} & \text{F} & \text{F} & \text{F} & \text{T} & \text{F} \\
\text{F} & \text{F} & \text{F} & \text{F} & \text{F} & \text{F} & \text{F} & \text{F} \\
\hline
\end{array}
\]
Columns 5 and 8 are identical.
\\
\textbf{Solution (b):}
\[
\begin{array}{|c|c|c|c|c|c|c|}
\hline
P & Q & P \lor Q & \sim(P \lor Q) & \sim P & \sim Q & \sim P \land \sim Q \\
\hline
\text{T} & \text{T} & \text{T} & \text{F} & \text{F} & \text{F} & \text{F} \\
\text{T} & \text{F} & \text{T} & \text{F} & \text{F} & \text{T} & \text{F} \\
\text{F} & \text{T} & \text{T} & \text{F} & \text{T} & \text{F} & \text{F} \\
\text{F} & \text{F} & \text{F} & \text{T} & \text{T} & \text{T} & \text{T} \\
\hline
\end{array}
\]
Columns 4 and 7 are identical.

\vspace{1em}
\textbf{Problem 2.61} Write negations of the open sentences: (a) Either $x = 0$ or $y = 0$. (b) The integers $a$ and $b$ are both even.
\\[0.5em]
\textbf{Solution:} By De Morgan's Laws:
\begin{itemize}
    \item \textbf{(a)} $x \neq 0$ and $y \neq 0$.
    \item \textbf{(b)} At least one of the integers $a$ or $b$ is odd (i.e., not even).
\end{itemize}

\vspace{1em}
\textbf{Problem 2.62} Implication: ``If $x$ and $y$ are even, then $xy$ is even.''
\\[0.5em]
\textbf{Solution:} 
Let $P$: ``$x$ and $y$ are even'', $Q$: ``$xy$ is even''.
\begin{itemize}
    \item \textbf{(a) Using `only if':} $x$ and $y$ are even only if $xy$ is even.
    \item \textbf{(b) Converse ($Q \Rightarrow P$):} If $xy$ is even, then $x$ and $y$ are even. (Note: mathematically false, but correctly formulated).
    \item \textbf{(c) Disjunction ($\sim P \lor Q$):} Either at least one of $x$ or $y$ is odd, or $xy$ is even.
    \item \textbf{(d) Negation as conjunction ($P \land \sim Q$):} $x$ and $y$ are even, and $xy$ is odd.
\end{itemize}

\vspace{1em}
\textbf{Problem 2.63} $P(x): x^2 = 2$ and $Q(x): x = \sqrt{2}$. State the negation of $P \Leftrightarrow Q$ in words.
\\[0.5em]
\textbf{Solution:} The negation of $P \Leftrightarrow Q$ is $(P \land \sim Q) \lor (\sim P \land Q)$.
\textbf{In words:} ``Either $x^2 = 2$ and $x \neq \sqrt{2}$, or $x^2 \neq 2$ and $x = \sqrt{2}$.''

\vspace{1em}
\textbf{Problem 2.64} Show that $[(P \lor Q) \land \sim(P \land Q)] \equiv \sim(P \Leftrightarrow Q)$.
\\[0.5em]
\textbf{Solution:}
\[
\begin{array}{|c|c|c|c|c|c|c|c|}
\hline
P & Q & P \lor Q & P \land Q & \sim(P \land Q) & (P \lor Q) \land \sim(P \land Q) & P \Leftrightarrow Q & \sim(P \Leftrightarrow Q) \\
\hline
\text{T} & \text{T} & \text{T} & \text{T} & \text{F} & \text{F} & \text{T} & \text{F} \\
\text{T} & \text{F} & \text{T} & \text{F} & \text{T} & \text{T} & \text{F} & \text{T} \\
\text{F} & \text{T} & \text{T} & \text{F} & \text{T} & \text{T} & \text{F} & \text{T} \\
\text{F} & \text{F} & \text{F} & \text{F} & \text{T} & \text{F} & \text{T} & \text{F} \\
\hline
\end{array}
\]
Columns 6 and 8 match. This is the definition of the ``Exclusive OR'' (XOR).

\newpage
\section*{Part 3: In-Depth Analysis of Problems and Techniques}

\subsection*{A) Problem Types and General Approach}
\begin{itemize}
    \item \textbf{Truth Table Verification:} Problems requiring you to definitively prove a tautology, contradiction, or equivalence (e.g., 2.46, 2.53, 2.60). \textit{Strategy:} Identify the number of unique variables, set up $2^n$ rows, evaluate inside parentheses first (PEMDAS for logic), and compare the final column against the desired outcome (all T's, all F's, or column-matching).
    \item \textbf{Algebraic Manipulation:} Instead of truth tables, reducing logical statements using laws (e.g., 2.51, 2.57). \textit{Strategy:} Immediately convert implications into disjunctions ($P \Rightarrow Q \equiv \sim P \lor Q$), distribute negations using De Morgan's, and look for known tautologies like $P \lor \sim P$ to collapse complex terms.
    \item \textbf{Translation (Math to English):} Applying logical rules to conversational sentences (e.g., 2.62, 2.63). \textit{Strategy:} Assign rigid variables ($P$, $Q$) to atomic clauses. Formulate the required logical construct symbolically, \textit{then} translate back to English. This prevents cognitive bias.
\end{itemize}

\subsection*{B) Key Algebraic Tricks and Manipulations}
\begin{itemize}
    \item \textbf{The Implication Elimination Trick:} Used heavily in 2.51 and 2.57. Because implications are hard to distribute over, changing $A \Rightarrow B$ to $\sim A \lor B$ unlocks the use of commutative and associative properties.
    \item \textbf{Contradiction Absorption:} Seen in 2.56. When a sub-clause evaluates to a known contradiction (like $P \land \sim P$), we instantly substitute it with $\text{F}$. Evaluating $\sim Q \Rightarrow \text{F}$ is vastly faster than expanding.
    \item \textbf{The Pigeonhole Argument:} Used in 2.59. A beautiful conceptual shortcut. Instead of writing out exhaustive tables for 5 undefined functions, realizing the limited ``state space'' (4 possible row outputs) forces overlaps. 
\end{itemize}


\section*{Part 4: ``Cheatsheet'' and Tips for Success}

\begin{itemize}
    \item \textbf{Implication Rule:} $P \Rightarrow Q$ is False in \textbf{only one scenario}: When Hypothesis ($P$) is T and Conclusion ($Q$) is F. In all other cases, it is T.
    \item \textbf{Vacuous Truth:} If $P$ is False, $P \Rightarrow Q$ is automatically True, regardless of $Q$. 
    \item \textbf{Inverse vs Converse:} Do not confuse them! The inverse of $P \Rightarrow Q$ is $\sim P \Rightarrow \sim Q$. The converse is $Q \Rightarrow P$. Both are logically equivalent to each other, but \textbf{not} to the original implication.
    \item \textbf{The Contrapositive:} $\sim Q \Rightarrow \sim P$ IS equivalent to the original $P \Rightarrow Q$.
    \item \textbf{Negating an Implication Trap:} The negation of $P \Rightarrow Q$ is \textbf{not} an implication. It is $P \land \sim Q$. (To break the promise, the condition was met, AND the outcome didn't happen).
    \item \textbf{Biconditional Negation:} $\sim(P \Leftrightarrow Q) \equiv (P \land \sim Q) \lor (\sim P \land Q)$. (They differ).
\end{itemize}

\section*{Part 5: Conceptual Synthesis and The ``Big Picture''}

\subsection*{A) Thematic Connections}
The core theme of Sections 2.7 to 2.9 is \textbf{structural invariance}. Just as $2x + 3x$ and $5x$ represent the same geometric line in algebra, logically equivalent statements represent the identical ``truth landscape.'' We are discovering that complex arguments can be reduced to simpler primitives, a theme that mirrors simplifying integrals in Calculus or diagonalizing matrices in Linear Algebra.

\subsection*{B) Forward and Backward Links}
Backward, this builds on Set Theory. Logical AND ($\land$) maps perfectly to Set Intersection ($\cap$), OR ($\lor$) to Union ($\cup$), and NOT ($\sim$) to Complements. De Morgan's laws govern both domains equivalently.
Forward, mastering logical equivalence is the undisputed bedrock of \textbf{Proof Writing}. When you encounter a theorem $P \Rightarrow Q$ that is too difficult to prove directly, you will fluidly switch to proving its contrapositive $\sim Q \Rightarrow \sim P$. When you write a ``Proof by Contradiction,'' you are actively employing the rule that assuming $P \land \sim Q$ leads to a contradiction (all F's), proving the original implication.

\newpage
\section*{Part 6: Real-World Application and Modeling (Finance Focus)}

\subsection*{A) Concrete Scenarios}
\begin{enumerate}
    \item \textbf{Barrier Options Pricing:} In quantitative finance, a "Down-and-In" call option only yields a payoff if the asset price breaches a specific barrier during the option's life AND the final asset price exceeds the strike price. A quant developer must code logical gates into Monte Carlo simulations to check these exact conditions.
    \item \textbf{Algorithmic Trading Executions:} A quantitative trading algorithm fires trades based on complex logical triggers. To optimize execution speed (measured in microseconds), quants use Boolean algebra to reduce the number of logical checks required before sending a buy order to an exchange.
    \item \textbf{Regulatory Stress Testing:} Under Basel III regulations, banks must halt stock buybacks if Capital Tier 1 drops below a threshold ($P$) AND economic volatility spikes ($Q$). Regulators analyze the negation of these policies to identify loopholes.
\end{enumerate}

\subsection*{B) Model Problem Setup: Algorithmic Logic Optimization}
\textbf{The Scenario:} A high-frequency trading desk has an algorithm to short an asset. The logic states: ``Execute short order ($S$) if the RSI is overbought ($R$) OR if the moving average crosses down ($M$), but DO NOT execute if liquidity is low ($L$).'' 
\\
\textbf{The Model:}
The initial naive logic built by junior developers is:
\[ S \Leftrightarrow ((R \lor M) \land \sim L) \]
However, checking $R$ and $M$ requires querying a database, which adds latency. The quant notices that if liquidity is low ($L$), the trade is blocked anyway. By applying distributive and implication logic, the quant reformulates the check:
The system should immediately abort if $L$ is True. We evaluate the equivalent condition for \textit{not} trading:
\[ \sim S \Leftrightarrow \sim((R \lor M) \land \sim L) \]
\[ \sim S \Leftrightarrow \sim(R \lor M) \lor L \]
\[ \sim S \Leftrightarrow (\sim R \land \sim M) \lor L \]
This tells the engine: "Check $L$ first. If True, abort. Then check $\sim R \land \sim M$. If True, abort." This mathematically identical model skips costly database calls, saving microseconds and increasing profitability.

\section*{Part 7: Common Variations and Untested Concepts}

\textbf{Concept: Algebraic Proofs of Equivalence without Truth Tables}\\
In standard curricula, exams frequently mandate: \textit{``Prove the statements are logically equivalent without using a truth table.''} This tests your ability to manipulate the symbolic algebra.
\\
\textbf{Example:} Prove $(P \land Q) \Rightarrow R \equiv P \Rightarrow (Q \Rightarrow R)$.
\\
\textbf{Solution:}
\begin{align*}
    (P \land Q) \Rightarrow R &\equiv \sim(P \land Q) \lor R \quad &\text{(Implication to Disjunction)} \\
    &\equiv (\sim P \lor \sim Q) \lor R \quad &\text{(De Morgan's Law)} \\
    &\equiv \sim P \lor (\sim Q \lor R) \quad &\text{(Associative Law)} \\
    &\equiv \sim P \lor (Q \Rightarrow R) \quad &\text{(Disjunction back to Implication)} \\
    &\equiv P \Rightarrow (Q \Rightarrow R) \quad &\text{(Disjunction back to Implication)}
\end{align*}
This elegantly demonstrates the equivalence without drawing an 8-row table.

\newpage
\section*{Part 8: Advanced Diagnostic Testing: ``Find the Flaw''}

\textit{Instructions: Below are 5 student solutions containing a subtle, critical error. Your task as a grader is to find the error, explain it in one sentence, and provide the correct steps. (The Answer Key is provided below for self-verification).}

\vspace{1em}
\hrule
\vspace{1em}
\textbf{Problem 1:} State the negation of the implication $P \Rightarrow Q$. \\
\textbf{Flawed Solution:} To negate an implication, we apply the negation to both sides. Therefore, $\sim(P \Rightarrow Q) \equiv \sim P \Rightarrow \sim Q$.

\vspace{1em}
\textbf{Problem 2:} Evaluate if $P \Rightarrow (P \land Q)$ is a tautology. \\
\textbf{Flawed Solution:} 
If $P$ is F, $P \land Q$ is F. The implication becomes $\text{F} \Rightarrow \text{F}$, which evaluates to F. Since it evaluates to False, it is not a tautology. 

\vspace{1em}
\textbf{Problem 3:} Simplify $\sim(P \Leftrightarrow Q)$. \\
\textbf{Flawed Solution:} The negation of a biconditional means both statements are no longer equivalent. Therefore, $\sim(P \Leftrightarrow Q) \equiv \sim P \Leftrightarrow \sim Q$. 

\vspace{1em}
\textbf{Problem 4:} Use the distributive law to expand $P \lor (Q \land R)$. \\
\textbf{Flawed Solution:} Distributing the OR operator yields $P \lor (Q \land R) \equiv (P \lor Q) \land R$. 

\vspace{1em}
\textbf{Problem 5:} Evaluate the truth value of $(P \land \sim P) \Rightarrow Q$. \\
\textbf{Flawed Solution:} $P \land \sim P$ is a contradiction (always False). Therefore, the premise is False. An implication with a false premise is invalid, meaning the statement evaluates to False. 

\vspace{1em}
\hrule
\vspace{1em}

\subsection*{Grader's Answer Key}

\textbf{Correction 1:} 
\textcolor{red}{Error:} The student confused the negation of an implication with its inverse. \\
\textcolor{blue}{Correct Step:} The actual negation translates to ``the condition was met, but the promise was broken,'' which is mathematically represented as $P \land \sim Q$.

\textbf{Correction 2:} 
\textcolor{red}{Error:} The student incorrectly evaluated $\text{F} \Rightarrow \text{F}$ as False. \\
\textcolor{blue}{Correct Step:} An implication with a false hypothesis is vacuously true, so $\text{F} \Rightarrow \text{F}$ evaluates to True (T). However, it is still not a tautology because if $P=\text{T}$ and $Q=\text{F}$, $\text{T} \Rightarrow \text{F}$ is False.

\textbf{Correction 3:} 
\textcolor{red}{Error:} $\sim P \Leftrightarrow \sim Q$ is actually logically equivalent to the original $P \Leftrightarrow Q$, not its negation. \\
\textcolor{blue}{Correct Step:} The negation of a biconditional is achieved by negating only one side (e.g., $\sim P \Leftrightarrow Q$), or explicitly writing it as $(P \land \sim Q) \lor (\sim P \land Q)$ (Exclusive OR).

\textbf{Correction 4:} 
\textcolor{red}{Error:} The student failed to distribute the $P$ term to the $R$ term. \\
\textcolor{blue}{Correct Step:} The proper distributive expansion requires pairing $P$ with both variables inside the parentheses: $(P \lor Q) \land (P \lor R)$.

\textbf{Correction 5:} 
\textcolor{red}{Error:} The student confused ``invalid reasoning'' with the truth value of a material implication; a false premise actually makes the implication vacuously True, not False. \\
\textcolor{blue}{Correct Step:} Since $P \land \sim P \equiv \text{F}$, the statement is $\text{F} \Rightarrow Q$. By the definition of material implication, this evaluates to True (T) regardless of $Q$.

\end{document}